\chapter{Pain relief (labour)}

\section*{Definition and Overview}
Pain relief during labour, or obstetric analgesia, refers to the spectrum of pharmacological and non-pharmacological interventions utilised to manage the pain associated with uterine contractions, cervical dilation, and fetal descent during childbirth. Labour pain is a complex, subjective, and multidimensional physiological response that carries significant clinical and psychological implications. The primary goal of labour pain management is to alleviate suffering, reduce maternal stress and exhaustion, and support maternal autonomy, while ensuring the safety of both the mother and the fetus. Management options range from non-invasive, self-administered techniques to advanced regional anaesthetic procedures, requiring an individualised approach tailored to the birthing person's preferences, medical history, and clinical progress.

\section*{Epidemiology}
Labour pain is a universal experience, although its perceived intensity and the subsequent utilization of pain relief options vary widely across demographics, cultural settings, and obstetric practices. In Australia, approximately 75\% of women giving birth utilize some form of pharmacological pain relief during labour. Nitrous oxide and oxygen (gas and air) is the most frequently used pharmacological intervention, employed in approximately 50\% of labours. Epidural analgesia is administered in approximately 30\% to 40\% of births, with higher rates observed in primiparous women, induced or augmented labours, and instrumental deliveries. Non-pharmacological pain relief methods, including water immersion, active birth positions, and continuous support, are widely adopted globally and are particularly prevalent in midwifery-led models of care.

\section*{Aetiology and Pathophysiology}
The painful stimuli experienced during labour are physiological and transition through distinct neuroanatomical pathways as labour progresses from the first to the second stage.

\begin{itemize}
    \item \textbf{Visceral pain (first stage of labour):} Driven primarily by uterine contractions, cervical effacement, and cervical dilation. Nociceptive signals from the uterus and cervix travel via visceral afferent A-delta and C fibres. These fibres accompany sympathetic nerves to enter the spinal cord at the T10 to L1 dermatomal levels. The pain is typically diffuse, dull, poorly localised, and referred to the lower abdomen, sacrum, and lower back.
    \item \textbf{Somatic pain (second stage of labour):} Arises from the mechanical stretching and compression of the vagina, perineum, pelvic floor, and surrounding skeletal structures as the fetus descends. Nociceptive signals are transmitted via somatic sensory fibres in the pudendal nerve, entering the spinal cord at the S2 to S4 dermatomal levels. This pain is sharp, intense, localised, and concentrated in the pelvic outlet and perineal region.
    \item \textbf{Uterine ischaemia:} During contraction, high intrauterine pressure compromises myometrial perfusion, leading to transient ischaemia and the release of inflammatory mediators (such as bradykinin, serotonin, and prostaglandins) that sensitise local nociceptors.
    \item \textbf{Psychological and cognitive factors:} Anxiety, fear, lack of preparation, and isolation can increase sympathetic nervous system activity and catecholamine release. This response causes maternal vasoconstriction, reduces uterine blood flow, increases uterine muscle tension, and lowers the pain threshold.
    \item \textbf{Obstetric variables:} Fetal malposition (such as persistent occiput posterior), cephalopelvic disproportion, and the administration of synthetic oxytocin (syntocinon) for induction or augmentation of labour significantly increase the intensity of nociceptive signals.
\end{itemize}

\section*{Clinical Features}
The clinical presentation of labour pain is progressive, dynamic, and correlates with the physiological phases of uterine activity and fetal descent.

\begin{itemize}
    \item \textbf{Rhythmic abdominal and pelvic pain:} Intermittent pain that correlates directly with uterine contractions, starting gradually, peaking in intensity, and then subsiding to a pain-free interval.
    \item \textbf{Referred lower back pain:} Often continuous and severe, particularly when the fetus is in an occiput posterior position, causing direct pressure on the maternal sacrum.
    \item \textbf{Perineal pressure and rectal urge:} Experienced during the second stage of labour, characterised by a deep, somatic pressure or burning sensation in the perineum, often accompanied by an involuntary urge to push.
    \item \textbf{Autonomic stress responses:} Fluctuations in maternal vital signs during contractions, including transient tachycardia, moderate hypertension, tachypnoea, and diaphoresis.
    \item \textbf{Psychomotor and emotional indicators:} Restlessness, vocalisation, facial grimacing, intense focus, withdrawal during contractions, and exhaustion between contractions.
\end{itemize}

\section*{Differential Diagnosis}
Although pain during labour is expected and physiological, it is critical to differentiate normal labour pain from acute, pathological conditions that require immediate obstetric or surgical intervention.

\begin{itemize}
    \item \textbf{Placental abruption:} Characterised by sudden, severe, and constant abdominal or back pain, uterine rigidity (woody feel) or hypertonus, and variable vaginal bleeding, which contrasts with the intermittent pain of physiological labour.
    \item \textbf{Uterine rupture:} Typically presents with a sudden, tearing abdominal pain, a cessation of contractions, loss of fetal station, maternal hypovolaemic shock, and acute fetal distress, occurring most frequently in patients with a history of prior Caesarean section.
    \item \textbf{Chorioamnionitis:} An intra-amniotic infection presenting with persistent abdominal tenderness between contractions, maternal pyrexia (temperature exceeding 38$^\circ$C), maternal and fetal tachycardia, and foul-smelling amniotic fluid.
    \item \textbf{Non-obstetric acute abdomen:} Pathologies such as appendicitis, cholecystitis, pancreatitis, or renal colic, which present with localised abdominal tenderness, guarding, rigidity, fever, or vomiting that do not correlate with uterine contractions.
\end{itemize}

\section*{When to Seek Medical Care}
Pregnant women should contact their midwife, obstetrician, or maternity hospital if they experience regular, painful contractions, a rupture of membranes (suspected leaking of amniotic fluid), vaginal bleeding, or a noticeable reduction in fetal movements.

\textbf{Call 000 (or local emergency services) for:}
\begin{itemize}
    \item Sudden, severe, and continuous abdominal pain that does not ease between contractions.
    \item Moderate to heavy, bright red vaginal bleeding.
    \item Loss of consciousness, maternal seizures (indicating eclampsia), or sudden severe shortness of breath.
    \item Feeling a loop of the umbilical cord in the vagina or vulva (umbilical cord prolapse) after the membranes rupture.
    \item An inability to cope with pain accompanied by maternal confusion, severe agitation, or physical collapse.
\end{itemize}

\section*{Diagnosis and Investigations}
The diagnosis of labour and the evaluation of associated pain require a comprehensive clinical assessment of the maternal-fetal unit.

\begin{itemize}
    \item \textbf{Clinical history and physical examination:} Assessment of contraction frequency, duration, and intensity via palpation, along with monitoring maternal vital signs (blood pressure, heart rate, temperature, and respiratory rate).
    \item \textbf{Digital vaginal examination:} Evaluation of cervical dilation (in centimetres), effacement (percentage), consistency, position, membrane status, and fetal station and presentation.
    \item \textbf{Cardiotocography (CTG):} Continuous or intermittent monitoring of the fetal heart rate and uterine contractions to assess fetal well-being and identify decelerations or loss of variability.
    \item \textbf{Maternal haematology:} A full blood count and group and screen are performed, particularly to check platelet count prior to regional anaesthesia (to exclude thrombocytopaenia) and to prepare for potential postpartum haemorrhage.
    \item \textbf{Urinalysis:} Screening for proteinuria (to assess for pre-eclampsia), ketones (to assess nutritional status and hydration), and leukocytes or nitrites (indicating urinary tract infection).
\end{itemize}

\section*{Management}
Management of labour pain involves a multidisciplinary, stepwise approach utilising both pharmacological and non-pharmacological strategies based on maternal preference and clinical circumstances.

\begin{itemize}
    \item \textbf{Epidural analgesia:} The most effective pharmacological option, involving the insertion of a catheter into the lumbar epidural space to administer a continuous infusion or patient-controlled boluses of local anaesthetics (e.g., ropivacaine, bupivacaine) and lipophilic opioids (e.g., fentanyl). It provides targeted sensory blockade of dermatomes T10 to S5.
    \item \textbf{Nitrous oxide and oxygen (Entonox):} Inhaled gas mixture (50\% nitrous oxide and 50\% oxygen) self-administered via a mouthpiece or mask. It is rapidly acting, cleared quickly through the lungs, and provides moderate analgesia and anxiolysis without significant neonatal accumulation.
    \item \textbf{Systemic opioids:} Intramuscular or intravenous administration of opioids, such as pethidine, morphine, or fentanyl (often delivered via Patient-Controlled Analgesia [PCA] pumps). These provide systemic analgesia, though they can cause maternal sedation, nausea, and neonatal respiratory depression if administered close to delivery.
    \item \textbf{Sterile water injections:} Subcutaneous or intradermal injections of sterile water (usually four injections of 0.1 mL to 0.2 mL) into the skin overlying the sacrum (Michaelis' rhomboid). This technique provides relief from severe back labour pain through the gate control theory of pain.
    \item \textbf{Non-pharmacological strategies:} Warm water immersion (labour pool), Transcutaneous Electrical Nerve Stimulation (TENS), massage, acupressure, heat packs, active birth positions, relaxation techniques, and continuous professional midwifery or doula support.
\end{itemize}

\section*{Complications}
The interventions used for pain relief in labour, particularly pharmacological options, carry specific risks and complications.

\begin{itemize}
    \item \textbf{Epidural-induced hypotension:} Maternal sympathetic blockade leading to peripheral vasodilation, which may reduce placental perfusion and cause transient fetal heart rate decelerations.
    \item \textbf{Post-dural puncture headache (PDPH):} A severe, postural headache occurring in 1\% of epidural insertions due to accidental puncture of the dura mater, causing cerebrospinal fluid leakage.
    \item \textbf{Prolonged second stage of labour:} Regional analgesia can reduce maternal expulsive efforts, increasing the duration of the second stage and slightly raising the rates of instrumental vaginal delivery (forceps or vacuum).
    \item \textbf{Maternal pyrexia:} Epidurals are associated with an unexplained elevation in maternal temperature, which can mimic infection and lead to neonatal investigations.
    \item \textbf{Neonatal respiratory depression:} Systemic opioids can cross the placenta, potentially causing neonatal sedation, low Apgar scores, and respiratory depression requiring neonatal resuscitation or naloxone administration.
    \item \textbf{Urinary retention:} Due to sensory and motor blockade of the bladder detrusor muscle, necessitating intermittent catheterisation.
\end{itemize}

\section*{Prognosis and Outcomes}
The vast majority of women achieve satisfactory pain relief, allowing them to navigate the physical challenges of childbirth successfully. The choice of pain relief, including epidural analgesia, does not increase the risk of emergency Caesarean section. While regional anaesthesia is associated with a minor increase in the rate of instrumental delivery, long-term complications are extremely rare. The overall birth experience and maternal psychological outcome depend not only on the efficacy of pain relief but also on the quality of supportive care, communication, and shared decision-making during the labour process.

\section*{Prevention and Self-Care}
Preparation and self-care strategies during pregnancy can assist women in managing labour pain and preparing for birth.

\begin{itemize}
    \item \textbf{Antenatal education:} Attending classes to understand the physiology of labour, options for pain relief, and birth preferences.
    \item \textbf{Formulating a flexible birth plan:} Documenting preferences for pain management, birth companions, and interventions, while remaining open to changes as labour progresses.
    \item \textbf{Active labour preparation:} Practising breathing techniques, mindfulness, meditation, and active birth positions (such as upright, kneeling, or squatting) to facilitate fetal positioning and comfort.
    \item \textbf{Physical conditioning:} Engaging in regular, low-impact exercise during pregnancy, such as prenatal yoga, swimming, and walking, to improve physical endurance.
    \item \textbf{Arranging support:} Selecting a supportive partner or professional birth companion (doula) to provide continuous emotional and physical support during labour.
\end{itemize}

\section*{Sources and Further Reading}
The following reputable organisations provide further information on pain relief during labour:

\begin{itemize}
    \item Royal Australian and New Zealand College of Obstetricians and Gynaecologists (RANZCOG). \textit{Pain relief in labour.} https://ranzcog.edu.au/
    \item Healthdirect Australia. \textit{Pain relief during labour.} https://www.healthdirect.gov.au/
    \item National Institute for Health and Care Excellence (NICE). \textit{Intrapartum care for healthy women and babies.} https://www.nice.org.uk/
    \item American College of Obstetricians and Gynecologists (ACOG). \textit{Obstetric analgesia and anesthesia.} https://www.acog.org/
    \item Pregnancy, Birth and Baby. \textit{Pain relief during labour.} https://www.pregnancybirthbaby.org.au/
\end{itemize}
